\chapter{Institutions and Long-Run Prosperity}\label{ch:b3:institutions}

Two places with the same climate, the same language and the same history can end up
with incomes that differ by a factor of twenty-five. The difference that the evidence
points to most often is the rules under which people produce and exchange: who owns
what, whether contracts are kept, whether the ruler can take what others have made.
This chapter measures those rules, reads the natural experiments with the methods of
the first chapters of this volume, and sets out how strong the case is and where it
is weak.

\begin{openingchart}{GDP per person in China, 1950--2022 (line; the dashed line marks 1978), and Maddison's own estimates for Western and Eastern Europe (marks;
logarithmic scale, 2011 dollars).}{Maddison Project Database 2023, country and regional
data}
\begin{tikzpicture}
\begin{semilogyaxis}[omtime, width=0.7\linewidth, height=0.44\linewidth, xmin=1950,
  xmax=2022, ymin=600, ymax=60000, ylabel={Dollars},
  ytick={1000,10000}, yticklabels={1{,}000, 10{,}000},
  legend pos=north west, legend style={font=\footnotesize}]
  \addplot[omThm, thick, no marks] table[x=year, y=gdppc] {figdata/bachelor-3/32-institutions/out/china.csv};
  \addplot[black, mark=*, mark size=1.5pt] table[x=year, y=gdppc] {figdata/bachelor-3/32-institutions/out/europe_west.csv};
  \addplot[omDef, mark=square*, mark size=1.5pt] table[x=year, y=gdppc] {figdata/bachelor-3/32-institutions/out/europe_east.csv};
  \draw[black!50, dashed] (axis cs:1978,600) -- (axis cs:1978,60000);
  \legend{China, Western Europe, Eastern Europe}
\end{semilogyaxis}
\end{tikzpicture}
\end{openingchart}

% figdata: opening_stats
China's income per person grew 2.8~\% a year from 1950 to 1978, from 799 to 1{,}744 dollars, and 5.5~\% a year from
1978 to 2022, to 19{,}238. Eastern Europe's income was 57~\% of Western Europe's in 1950 and 40~\% in 1990, after four
decades of planning; by 2022 it was 50~\%.

\begin{recall}
\omterm{def:b3:institutions:institutions}{Institutions}, the two Koreas, the two Germanys and China's reforms appear in the
grades 6--9 and 10--12 volumes; growth accounting and misallocation in the Year 1
volume; prices as information in the Year 2 volume. \omterm{def:b3:instrumental-variables:instrument}{Instrumental variables},
difference in differences with \omterm{def:b3:difference-in-differences:synthetic-control}{synthetic control}, and regression discontinuity are
\cref{ch:b3:instrumental-variables,ch:b3:difference-in-differences,ch:b3:regression-discontinuity}'s.
\end{recall}

\begin{notation}[Symbols of this chapter]\label{not:b3:institutions:symbols}
$p$ is the probability that output is expropriated; $k$ investment; $A k^{a}$ output; $z$ an
instrument; $s$ a measure of property security; $\pi$ a first-stage slope.
\end{notation}

\section{Defining and measuring institutions}\label{sec:b3:institutions:defining}

\begin{definition}[Institutions]\label{def:b3:institutions:institutions}
\index{institutions}
\emph{Institutions} are the rules of the game in a society --- laws, \omterm{def:b3:institutions:property-rights}{property rights},
the constraints on rulers and the ways they are enforced --- as distinct from the
players' choices within them. In data they are measured by indices built from
surveys of experts, firms and households, or by recorded rules such as the share of
land with registered titles.
\end{definition}

\begin{definition}[Property rights; rule of law; expropriation risk]\label{def:b3:institutions:property-rights}
\index{property rights}\index{rule of law}\index{expropriation risk}
\emph{Property rights} are the rights to use an asset, to its income and to sell it.
The \emph{rule of law} is the extent to which rules bind everyone, rulers included,
and are enforced impartially, measured as an index. \emph{Expropriation risk} is the
probability that the state or others take an asset or its income.
\end{definition}

\begin{proposition}[Expropriation lowers investment]\label{prop:b3:institutions:holdup}
An investor chooses $k$ to maximise $(1 - p)A k^{a} - k$, with $0 < a < 1$: the output is taken with probability $p$ after the
investment is sunk. Then $k^* = (a(1 - p)A)^{1/(1 - a)}$, which falls with $p$, and output $A k^{*a}$ is proportional to $(1 - p)^{a/(1 - a)}$.
\end{proposition}

\begin{proof}
The objective is strictly concave in $k$. The first-order condition $(1 - p)aA k^{a - 1} = 1$ gives $k^*$, and $dk^*/dp = -k^*/((1 - a)(1 - p)) < 0$.
Substituting, $A k^{*a} = A(aA)^{a/(1 - a)}(1 - p)^{a/(1 - a)}$.
\end{proof}

With $a = 1/3$, output is proportional to $(1 - p)^{1/2}$: a risk of 50~\% lowers output by 29~\%, and more once the
investment is in skills, roads and firms that a short horizon also discourages.

\begin{measure}[Measuring institutions]
Indices of the \omterm{def:b3:institutions:property-rights}{rule of law} or of \omterm{def:b3:institutions:property-rights}{expropriation risk} are built from experts' and
firms' assessments, so they mix the rules with the outcomes: a rich country's
courts look reliable partly because the country is rich. Most indicators used to
show that \omterm{def:b3:institutions:institutions}{institutions} cause growth measure outcomes, not constraints, and some of
the instruments used are flawed. % ledger: b3in:glls
Recorded rules --- titled land, the date of an independent court --- are closer to
the definition but cover less.
\end{measure}

\section{The natural experiments}\label{sec:b3:institutions:experiments}

\begin{omchart}{GDP per person in South and North Korea, 1990--2022 (logarithmic scale, 2011
dollars; the North Korea series starts in 1990).}{Maddison Project Database 2023}
\begin{tikzpicture}
\begin{semilogyaxis}[omtime, width=0.62\linewidth, height=0.4\linewidth, xmin=1990,
  xmax=2022, ymin=1000, ymax=50000, ylabel={Dollars},
  ytick={1000,10000}, yticklabels={1{,}000, 10{,}000},
  legend style={font=\footnotesize, at={(0.97,0.5)}, anchor=east}]
  \addplot[omDef, thick] table[x=year, y=gdppc] {figdata/bachelor-3/32-institutions/out/korea_south.csv};
  \addplot[omThm, thick] table[x=year, y=gdppc] {figdata/bachelor-3/32-institutions/out/korea_north.csv};
  \legend{South Korea, North Korea}
\end{semilogyaxis}
\end{tikzpicture}
\end{omchart}
% figdata: koreas_stats

The peninsula was divided in 1945 into two states with opposite rules: private
property and markets in the South, collective property and central planning in the
North. In 1990 the South's income per person was 5.6 times the North's; in 2022 it was 26.3 times, the
North at 1{,}569 dollars. The division is close to an experiment: the border was drawn by
outside powers along a line of latitude, not along any prior difference in the
economy. If anything, the starting point favoured the North, which held 80~\% of the
peninsula's coal and minerals and most of the heavy machinery left from colonial rule,
while most of the South's people worked small farms. % ledger: b3in:cfr.korea
It is not a clean experiment: the two states also differed in their allies, their
aid and the wars they prepared for, and they are two units, which allows no
standard error.

Germany's division is the second case. Its use as an experiment goes beyond the
comparison of the two states: after division, West German cities close to the new
border saw their population growth fall well behind that of other West German
cities, as a model of market access predicts --- the border cut them off from their
hinterland. % ledger: b3in:rs
The comparison of the two states is the less clean one, since the East also paid
reparations and lost workers until the border was closed.

\begin{record}[China's household responsibility system, 1978--1984]
From 1978 Chinese farm households were allowed to keep what they produced beyond a
quota, and the collectives were dissolved. A provincial panel finds that
decollectivisation raised total factor productivity and accounts for about half of
the growth of agricultural output in 1978--84; higher state procurement prices helped too,
mainly by raising the use of inputs, while other market reforms mattered little. % ledger: b3in:lin
The timing differed across provinces, which is what lets the panel separate the
reform from the common trend (\cref{ch:b3:difference-in-differences}).
\end{record}

\begin{model}{Institutions as the cause of income differences}
\assumes Secure \omterm{def:b3:institutions:property-rights}{property rights} and constraints on rulers raise the return to
investment in capital, skills and firms
(\cref{prop:b3:institutions:holdup}).
\predicts Large income gaps between neighbours that share geography and culture but
not rules; income responding to changes of rules.
\tested The two Koreas and the two Germanys; China after 1978; instruments for
\omterm{def:b3:institutions:institutions}{institutions} across former colonies (next section); a regression discontinuity on
a colonial labour system in Peru. % ledger: b3in:dell
\breaks The measures of \omterm{def:b3:institutions:institutions}{institutions} are partly outcomes; \omterm{def:b3:human-capital:human-capital}{human capital} may come
first, with poor countries growing under good policies, often under dictators, and
improving their political \omterm{def:b3:institutions:institutions}{institutions} afterwards. % ledger: b3in:glls
\end{model}

\section{Instruments for institutions}\label{sec:b3:institutions:instruments}

Regressing income on an index of \omterm{def:b3:institutions:institutions}{institutions} across countries runs into all three
problems of \cref{ch:b3:instrumental-variables}: rich countries can afford better
\omterm{def:b3:institutions:institutions}{institutions}, the index is measured with error, and omitted factors drive both.

\begin{definition}[Settler-mortality instrument]\label{def:b3:institutions:settler-mortality}
\index{settler-mortality instrument}
The \emph{settler-mortality instrument} uses the death rates that European soldiers,
bishops and sailors faced in the colonies from the seventeenth to the nineteenth
centuries as an instrument for current \omterm{def:b3:institutions:institutions}{institutions}. Where Europeans could settle
they built \omterm{def:b3:institutions:institutions}{institutions} that protected property; where they died they set up
extractive ones; and the early \omterm{def:b3:institutions:institutions}{institutions} persisted. % ledger: b3in:ajr
\end{definition}

\begin{omfigure}{The causal chain behind the settler-mortality instrument. The exclusion
restriction is that mortality affects income today only through institutions (the
crossed arrow must be absent).}
\begin{tikzpicture}[node distance=2.6cm, every node/.style={font=\small}]
  \node[draw, rounded corners, align=center] (m) {Settler\\mortality};
  \node[draw, rounded corners, align=center, right=of m] (s) {Settlement};
  \node[draw, rounded corners, align=center, right=of s] (i) {Early\\institutions};
  \node[draw, rounded corners, align=center, right=1.6cm of i] (c) {Current\\institutions};
  \node[draw, rounded corners, align=center, below=1.4cm of c] (y) {Income\\today};
  \draw[-{Stealth}, thick] (m) -- (s);
  \draw[-{Stealth}, thick] (s) -- (i);
  \draw[-{Stealth}, thick] (i) -- (c);
  \draw[-{Stealth}, thick] (c) -- (y);
  \draw[-{Stealth}, dashed, omThm] (m) |- node[pos=0.75, fill=white] {$\times$ health, disease today} (y);
\end{tikzpicture}
\end{omfigure}

\begin{proposition}[Two-stage least squares with one instrument]\label{prop:b3:institutions:wald}
With one \omterm{def:b3:instrumental-variables:endogenous}{endogenous regressor} $s$ and one instrument $z$, the \omterm{def:b3:instrumental-variables:tsls}{two-stage least squares}
coefficient equals the slope of the \omterm{def:b3:instrumental-variables:tsls}{reduced form} (income on $z$) divided by the slope
of the \omterm{def:b3:instrumental-variables:tsls}{first stage} ($s$ on $z$).
\end{proposition}

\begin{proof}
The \omterm{def:b3:instrumental-variables:tsls}{first stage} gives $\hat s = c + \pi z$, with $\pi = \operatorname{cov}(z, s)/\operatorname{var}(z)$. The second stage regresses $\ln y$ on $\hat s$:
its slope is $\operatorname{cov}(\hat s, \ln y)/\operatorname{var}(\hat s) = \pi\operatorname{cov}(z, \ln y)/(\pi^2\operatorname{var}(z)) = (\operatorname{cov}(z, \ln y)/\operatorname{var}(z))/\pi$, the reduced-form slope over the first-stage slope.
\end{proof}

The ratio makes the weak point plain: when the \omterm{def:b3:instrumental-variables:tsls}{first stage} is small and imprecise,
the estimate is a ratio with a denominator near zero
(\cref{ch:b3:instrumental-variables}).

\begin{debate}[Geography, institutions or human capital?]
\textbf{\omterm{def:b3:institutions:institutions}{Institutions} first.} With settler mortality as the instrument, differences in
\omterm{def:b3:institutions:institutions}{institutions} explain about three-quarters of the differences in income per person across
former colonies, and once \omterm{def:b3:institutions:institutions}{institutions} are controlled for, being in Africa or far
from the equator does not lower income (\emph{Grade 3}). % ledger: b3in:ajr
\emph{Weakness:} the \omterm{def:b3:instrumental-variables:instrument}{exclusion restriction} --- mortality then may affect income now
through disease and health as well.

\textbf{The data behind the instrument.} In 36 of the 64 countries of the sample, the mortality rate is
taken from another country, and rates from labourers, bishops and soldiers on
campaign are combined; with these issues controlled for, the link between mortality
and \omterm{def:b3:institutions:property-rights}{expropriation risk} is not robust and the instrumental-variable estimates become
unreliable, often with infinite confidence intervals. % ledger: b3in:albouy
\emph{Weakness:} a critique of one instrument, not evidence against the effect.

\textbf{Geography directly.} Malaria transmission, set largely by ecology, lowers
income per person even after the quality of \omterm{def:b3:institutions:institutions}{institutions} is controlled for (\emph{Grade 3}). % ledger: b3in:sachs
\emph{Weakness:} malaria is itself shaped by income and public health.

\textbf{\omterm{def:b3:human-capital:human-capital}{Human capital} first.} Schooling may be the more basic source of growth,
with \omterm{def:b3:institutions:institutions}{institutions} improving after it (\emph{Grade 3}). % ledger: b3in:glls
\emph{Weakness:} cross-country OLS, open to the reverse causation it criticises.

\textbf{Persistence, measured locally.} A regression discontinuity at the boundary of
the mita, the forced mining labour imposed in Peru and Bolivia from 1573 to 1812, finds household
consumption about 25~\% lower and child stunting about 6 points higher inside it today, through land
tenure and public goods: fewer large landowners historically, less schooling, fewer
roads (\emph{Grade 2}). % ledger: b3in:dell
\emph{Weakness:} one boundary in one country; it shows that \omterm{def:b3:institutions:institutions}{institutions} persist, not
how much of world inequality they explain.
\end{debate}

\section{Planned economies in the data}\label{sec:b3:institutions:planned}

\begin{definition}[Planned economy; soft budget constraint]\label{def:b3:institutions:planned}
\index{planned economy}\index{soft budget constraint}
A \emph{planned economy} allocates resources by administrative plan, with state
ownership of firms and prices set by planners. A firm faces a \emph{soft budget
constraint} when its losses will be covered by another institution, typically the
state; the higher the probability that losses are covered, the softer the constraint. % ledger: b3in:kornai
\end{definition}

The softness weakens firms' response to prices, lowers efficiency and can create
excess demand; it was extreme in the socialist economies and is present to a lesser
degree in mixed ones. % ledger: b3in:kornai
Exercise 11 shows the mechanism: a firm that expects to be bailed out puts in less
effort to avoid failure. The opening chart shows the result at the scale of a
continent: Eastern Europe fell from 57 to 40~\% of Western Europe's income between 1950 and 1990.

\section{The transitions of the 1990s}\label{sec:b3:institutions:transition}

\begin{definition}[Transition recession]\label{def:b3:institutions:transition}
\index{transition recession}
A \emph{transition recession} is the fall of measured output that followed the end of
central planning, as old links between firms broke before new markets formed.
\end{definition}

\begin{omchart}{GDP per person in four economies after central planning, as an index with
1989 = 100.}{Maddison Project Database 2023 (successor states' own series)}
\begin{tikzpicture}
\begin{axis}[omtime, width=0.68\linewidth, height=0.42\linewidth, xmin=1985,
  xmax=2022, ymin=40, ymax=380, ylabel={Index, 1989 = 100},
  legend pos=north west, legend style={font=\footnotesize},
  extra y ticks={100}, extra y tick labels={}, extra y tick style={grid=major}]
  \addplot[omDef, thick] table[x=year, y=idx] {figdata/bachelor-3/32-institutions/out/transition_POL.csv};
  \addplot[omDef, dashed, thick] table[x=year, y=idx] {figdata/bachelor-3/32-institutions/out/transition_CZE.csv};
  \addplot[omThm, thick] table[x=year, y=idx] {figdata/bachelor-3/32-institutions/out/transition_RUS.csv};
  \addplot[omThm, dashed, thick] table[x=year, y=idx] {figdata/bachelor-3/32-institutions/out/transition_UKR.csv};
  \legend{Poland, Czechia, Russia, Ukraine}
\end{axis}
\end{tikzpicture}
\end{omchart}
% figdata: transition_stats

Poland's income per person fell 16~\% to 1991 and was back above its 1989 level by 1995; Czechia's fell
9~\%. Russia's fell 34~\% and bottomed out only in 1996; Ukraine's fell 54~\%, to 1997, regained the 1989 level
in 2007, and in 2022 was at 77, after the war. Measured output overstates the fall where
planned output had no buyers and understates it where queues disappeared; but the
differences between countries are too large to be measurement alone. The countries
also differed in how fast they freed prices, how they privatised and how much their
courts and rulers changed, and they started from different places; no simple
comparison of two countries separates these.

\section{Crime, deterrence and the rule of law}\label{sec:b3:institutions:crime}

Secure property also means protection from theft. Whether police reduce crime is
hard to measure: cities with more crime hire more police, so a simple regression
can even show police going with more crime.

\begin{definition}[Deterrence elasticity of crime]\label{def:b3:institutions:deterrence}
\index{deterrence elasticity of crime}
The \emph{deterrence elasticity of crime} is the percentage change in crime for a
1~\% increase in the resources of enforcement, such as the number of police.
\end{definition}

Instruments and natural experiments try to break the simultaneity, and the first
attempt shows how fragile an instrument can be.

\begin{debate}[How much do police reduce crime?]
In large U.S. cities in 1970--1992, police forces grew disproportionately in mayoral and
governors' election years. With elections as the instrument, Levitt (1997) found that
police reduce crime in six of seven categories, each extra officer preventing eight to
ten serious crimes (\emph{Grade 3} as published: an instrument whose exclusion rests
on elections affecting crime only through police). % ledger: b3in:levitt
Replicating the study with the data and code its author supplied, McCrary (2002)
found a weighting error in the program: with the intended weights, every estimate is
statistically insignificant, even at the 10~\% level, and a better measure of election
years strengthens the \omterm{def:b3:instrumental-variables:tsls}{first stage} but leaves the estimates less precise still
(\emph{Grade 3}: the instrument moves police too little to measure their effect).
% ledger: b3in:mccrary
Levitt's reply (2002) turned to another instrument, the number of firefighters per
head, which tracks city budgets, and found police reducing both violent and property
crime, most clearly murder and robbery (\emph{Grade 3}: the instrument is valid only
if budgets for firefighters do not respond to crime). % ledger: b3in:levitt02
Over 1960--2010, Chalfin and McCrary argued that the larger problem is error in the
recorded size of police forces, which biases estimates towards zero; correcting for
it, they put the elasticity of the cost-weighted total of crimes at about $-0.5$, larger
for violent crime than for property crime (\emph{Grade 2--3}). % ledger: b3in:cm
After the July 1994 attack on the main Jewish centre in Buenos Aires, every Jewish
institution received police protection; car thefts fell sharply on the protected
blocks, with no appreciable change beyond them (\emph{Grade 1--2}: an allocation of
police unrelated to car theft, though the effect is local and may move theft
elsewhere). % ledger: b3in:dts
The local experiment shows that visible police deter where they stand; the size of
the city-wide elasticity, which is what a hiring decision needs, remains contested.
\end{debate}

\section{Exercises}\label{sec:b3:institutions:exercises}

\begin{exercise}[$\star$]\label{exo:b3:institutions:1}
From the opening chart's numbers, find China's growth rate before and after 1978, and how
many times faster it grew after.
\end{exercise}

\begin{exercise}[$\star$]\label{exo:b3:institutions:2}
The South's income per person in 2022 was 41{,}321 dollars and the North's 1{,}569. Find the ratio and the number of
doublings it represents.
\end{exercise}

\begin{exercise}[$\star$]\label{exo:b3:institutions:3}
In an invented sample, the slope of log income on log settler mortality is $-0.57$ and the
slope of an \omterm{def:b3:institutions:institutions}{institutions} index on log mortality is $-0.61$. Find the \omterm{def:b3:instrumental-variables:tsls}{two-stage least squares}
estimate.
\end{exercise}

\begin{exercise}[$\star$]\label{exo:b3:institutions:4}
What evidence would contradict the view that \omterm{def:b3:institutions:institutions}{institutions} caused most of the gap
between the Koreas?
\end{exercise}

\begin{exercise}[$\star\star$]\label{exo:b3:institutions:5}
In \cref{prop:b3:institutions:holdup} take $a = 1/2$ and $A = 2$. Find investment and output for $p = 0$ and $p = 0.3$.
\end{exercise}

\begin{exercise}[$\star\star$]\label{exo:b3:institutions:6}
Corrected mortality data weaken the \omterm{def:b3:instrumental-variables:tsls}{first stage} of exercise 3 to $-0.20$, with the \omterm{def:b3:instrumental-variables:tsls}{reduced form}
unchanged. What happens to the estimate and to its precision?
\end{exercise}

\begin{exercise}[$\star\star$]\label{exo:b3:institutions:7}
Compute the depth, the length and the recovery year of the \omterm{def:b3:institutions:transition}{transition recessions} in
Poland, Czechia, Russia and Ukraine.
\end{exercise}

\begin{exercise}[$\star\star$]\label{exo:b3:institutions:8}
In an invented city, car thefts per block and month are 0.09 on blocks with a protected
institution and on other blocks before an attack; after it, 0.03 on the protected blocks and
0.085 on the others. Find the difference in differences and say what pattern would show
that theft had only moved.
\end{exercise}

\begin{exercise}[$\star\star$]\label{exo:b3:institutions:9}
What pattern at the mita boundary would break the regression discontinuity?
\end{exercise}

\begin{exercise}[$\star\star\star$]\label{exo:b3:institutions:10}
Let $\ln y = \beta s + \gamma h + u$, where $h$ is health today, and let both $s$ and $h$ depend on mortality $z$ with
slopes $\pi_s$ and $\pi_h$. Find what the instrumental-variable estimate with $z$ converges to, and
its bias.
\end{exercise}

\begin{exercise}[$\star\star\star$]\label{exo:b3:institutions:11}
A firm's manager chooses effort $e$ in $[0, 1]$ at cost $e^2/2$; the project succeeds with
probability $e$ and pays $R$. On failure, the state covers the losses with probability $b$,
paying $B$. Find effort and show that it falls with $b$. Take $R = 0.8$, $B = 0.6$.
\end{exercise}

\begin{exercise}[$\star\star\star$]\label{exo:b3:institutions:12}
A city has 1{,}000 officers and 50{,}000 serious crimes a year, and considers hiring 100 more
at 80{,}000 dollars each a year. Take three values of the \omterm{def:b3:institutions:deterrence}{deterrence elasticity of crime} from
the debate: 0 (the corrected election estimates are not distinguishable from zero),
$-0.18$ (nine crimes prevented per officer, as in the original estimate) and $-0.5$. For
each, find the crimes prevented per officer and the social cost per crime above which
hiring pays. Which estimate decides the verdict, and what would you need to know
about the crimes prevented?
\end{exercise}

\section{Problem: Brenland and the Land across the River}\label{sec:b3:institutions:problem}

\begin{problem}[{Weekend problem --- Brenland and the land across the river}]\label{pb:b3:institutions:river}
Two invented provinces of Brenland face each other across the river Varn. They
shared one path of income until 1995, when North registered titles to land and firms;
South did not. Eight other provinces had no reform. A survey of 2025 covers 60
districts, 30 on each bank: the share of titled plots $s$, log income per person, and the
distance $z$ to the land-survey office of 1890, which titled plots near it first. % figdata: river_stats

\medskip\noindent\textbf{Part I --- the gap.}
\begin{enumerate}
\item Find income per person on each bank in 1980, 1995 and 2025.
\item Find the log gap in 1995 and in 2025.
\item Find each bank's growth rate before and after the reform.
\item Find the difference in differences of log income, 1995 to 2025.
\item What must hold before 1995 for it to measure the reform's effect? Does it?
\item Why might it still not measure the effect of the titles?
\end{enumerate}

\medskip\noindent\textbf{Part II --- a synthetic North.}
\begin{enumerate}[resume]
\item Build a synthetic North from the eight other provinces, matching 1980--1995. Give
the weights.
\item Find the root mean squared error of the fit before 1995.
\item Find the gap between North and its \omterm{def:b3:difference-in-differences:synthetic-control}{synthetic control} in 2025.
\item Compare it with the difference in differences.
\item Why is South left out of the donors?
\item What would a large gap before 1995 mean?
\end{enumerate}

\medskip\noindent\textbf{Part III --- titles and income across districts.}
\begin{enumerate}[resume]
\item Regress log income on $s$ by least squares.
\item Find the \omterm{def:b3:instrumental-variables:tsls}{first stage} of $s$ on $z$ and its $F$ statistic.
\item Find the \omterm{def:b3:instrumental-variables:tsls}{reduced form} of log income on $z$.
\item Find the \omterm{def:b3:instrumental-variables:tsls}{two-stage least squares} estimate.
\item Why is least squares higher?
\item Name one way distance to the survey office could affect income other than
through titles.
\item The banks differ by $\Delta s$ in mean titled share and by $\Delta\ln y$ in mean log income. Find the share of
the gap the titles explain, by 2SLS and by least squares.
\end{enumerate}

\medskip\noindent\textbf{Part IV --- the investment model.}
\begin{enumerate}[resume]
\item With \cref{prop:b3:institutions:holdup} and $a = 1/3$, write the output ratio of the banks in terms of
their \omterm{def:b3:institutions:property-rights}{expropriation risks}.
\item The titles explain a log gap of $\beta_{\mathrm{2SLS}}\Delta s$. With a risk of 5~\% in North, what risk in South
produces it?
\item Is that risk plausible for untitled land? What else would the titles change?
\item Why does the model's gap understate the effect of titles if titled land also
serves as collateral?
\item State the result: the share of the income gap explained by property security,
by 2SLS.
\end{enumerate}
\end{problem}
